\section{A unified object, and how to falsify these claims}
\label{sec:unified}

\subsection{What the three results describe together}

Placed side by side, the surviving results do not describe three patches but one
design, with each axis constrained by a different theorem:
\begin{itemize}
  \item \textbf{State.} A learned, bounded, gated recurrence
        (Theorem~\ref{thm:linattn}) --- but with a \emph{dense} transition if an
        HMM-like reading is wanted (Theorem~\ref{thm:noembed}), and with enough
        exact-recall capacity to clear Theorem~\ref{thm:recall}.
  \item \textbf{Time.} Serial steps that need not emit a token
        (Theorem~\ref{thm:sim}), but which \emph{accumulate} rather than
        overwrite (Corollary~\ref{cor:crossover}), with projection onto a learned
        codebook every $k$ steps (Open Problem~\ref{op:drift}).
  \item \textbf{Depth.} Early exit governed by tail energy
        (Theorem~\ref{thm:accept}), supplying both an adaptive halting signal and,
        on the tokens that do not halt, a free speculative draft
        (Theorem~\ref{thm:dominance}).
\end{itemize}
The depth axis serves two purposes at once: the signal that says ``this token is
easy, emit it now'' is the signal that produces a draft for the tokens that are
not. Adaptive computation time \cite{graves2016act} and self-speculation are the
same mechanism read in opposite directions. We prove nothing about the
combination and regard the experiments below as prior to any such attempt.

\subsection{Protocol}

Each experiment targets one theorem and carries a kill criterion fixed in
advance. Models are $\le$1B parameters; two experiments require no training.

\paragraph{E1: acceptance profile (Theorems~\ref{thm:accept},
\ref{thm:dominance}).} For every layer $\ell$ of Qwen3-0.6B, measure
$\alpha_\ell = 1 - \TV(p_\ell,p_L)$, tail energy $T_\ell$, and the bound
\eqref{eq:acceptbound}, then evaluate the speedup surface with the corrected cost
\eqref{eq:cnoreuse} including the head fraction $u \approx 0.26$. Report the
looseness of \eqref{eq:acceptbound} decomposed into its Cauchy--Schwarz and
softmax stages, and compare against Proposition~\ref{prop:rank} with $P$ fitted
to the post-norm tail. \textbf{Kill:} $\max_{\ell,\gamma} S \le 1$ both untuned
and after an early-exit auxiliary-loss adaptation.

\paragraph{E2: bandwidth (Theorem~\ref{thm:stream}).} Train $\sim$10M-parameter
models from scratch on $\mathrm{PROD}_{N,T}$ under a \emph{streaming} read
schedule, which Proposition~\ref{prop:confound} shows must be enforced rather
than assumed: the re-read escape is otherwise open, and because composition is
associative, recomputation costs only $\log_2 T$ depth, so a $4$-layer model
needs $T \gtrsim 2^{8}$ before depth alone closes it. Sweep the interface width
against $\log_2(N!)$ and check the predicted threshold at $N = 9$
(Corollary~\ref{cor:instantiate}). \textbf{Kill:} accuracy does not fall at
$b < \log_2(N!)$ once the read schedule is enforced.

\paragraph{E3: the HMM identity and the missing bridge
(Theorems~\ref{thm:hmm}, \ref{thm:noembed}).} Verify Theorem~\ref{thm:hmm}
numerically against brute-force path enumeration, per state and per timestep, with
a negative control that drops non-negativity. Then verify
Theorem~\ref{thm:noembed} directly: check that row-support is non-increasing
along a gated linear-attention recurrence and strictly increasing along the
forward algorithm, and that the $\Psi_o$ fail to commute. \textbf{Kill:} the
identity fails, or the support and commutator tests come out the other way.

\paragraph{E4: capacity (Theorem~\ref{thm:recall}).} Restrict Qwen3-0.6B's
attention to a sliding window of $w$ tokens as a bounded-state proxy and locate
the associative-recall break-point as a function of $w$. Note that the Fano form
of Theorem~\ref{thm:recall} predicts \emph{graceful} degradation past threshold
rather than a cliff. \textbf{Kill:} the break-point does not move with $w$, which
would falsify the proxy rather than the theorem.

\subsection{Limitations}

The cost model of Definition~\ref{def:cost} ignores attention's quadratic term
and KV traffic. Theorem~\ref{thm:tc0} is a statement about classes and does not
constrain any individual trained model (Remark~\ref{rem:tc0-care}).
Theorem~\ref{thm:sim} and Theorem~\ref{thm:tc0} assume incompatible precision
regimes. Theorem~\ref{thm:stream} requires a streaming read schedule that real
architectures do not obey, and Remark~\ref{rem:barrington} explains why removing
that hypothesis is not merely difficult but would separate $\mathrm{TC}^0$ from
$\mathrm{NC}^1$. Finally, no experiment in this paper has been run: the results
above are stated as predictions with kill criteria, and we report them as such.
